\documentclass[12pt]{article}
\usepackage[margin=1in]{geometry}
\usepackage{parskip}
\usepackage{booktabs}
\usepackage{hyperref}

\title{\textbf{Assignment 2 } \\
\large {}}
\author{}
\date{}

\begin{document}

\maketitle

\section*{Interpretation of Findings}

\subsection*{Estimated Marginal Productivity Gain}

The randomized controlled trial shows that the AI PDF extraction tool significantly increased clerk productivity. Clerks using the tool completed an average of 80.50 tasks per shift, compared to 69.95 for the control group, resulting in an Average Treatment Effect (ATE) of +10.55 tasks, or a 15.1\% increase. Because assignment to the AI tool was randomized, this improvement can be directly attributed to the tool rather than pre-existing differences between clerks. The balance test further confirms that both groups were statistically similar before the intervention, supporting the validity of the results.

\subsection*{Quality and the Quantity--Quality Trade-Off}

While the AI tool increased output, it did not improve quality. Clerks using the tool had a slightly higher error rate (2.69\%) compared to the control group (2.25\%), resulting in an ATE of +0.45\%. Although small, this suggests a quantity–quality trade-off: clerks processed more tasks but made slightly more errors. This may reflect over-reliance on automation or reduced attention due to faster processing speeds.

\subsection*{Implications for Firms Considering Broad AI Deployment}

These findings have important implications for firms considering these tools at scale. The productivity gains are clear and make a strong case for improving efficiency. However, the increase in error rates is a concern, especially in loan processing where mistakes can lead to financial and compliance issues.

To address this, firms should pair adoption with strong quality control measures, such as tracking error rates, conducting regular audits, and ensuring employees are properly trained on how to use the tool effectively. A phased rollout with ongoing performance monitoring is recommended before full implementation. The randomized controlled trial provides a reliable benchmark, and firms should consider running similar pilot tests to evaluate the impact before scaling.

\end{document}
